\documentclass[12pt]{article}
\usepackage[T1]{fontenc}
\usepackage[utf8]{inputenc}
\usepackage{lmodern}
\usepackage{textcomp}
\usepackage{enumitem}
\usepackage{geometry}
\geometry{margin=1in}
\begin{document}
\begin{center}
Answer Key - Biology - Graduate Level Assessment 11 \
\small (Answers are ordered exactly as in the question paper.)
\end{center}

\section*{Multiple Choice}
\begin{enumerate}[label=\arabic*.]
\item \textbf{Answer:} B
\\ \textit{Options:} A) parenchyma and companion cells; B) tracheids and vessel elements; C) parenchyma and collenchyma; D) sieve tube members and companion cells; E) collenchyma and vessel elements
\\ \textit{Note:} Tracheids and vessel elements are specialized xylem cells that facilitate the transport of water and minerals while also providing structural support to the plant, making them essential for both functions in vascular plants.
\item \textbf{Answer:} B
\\ \textit{Options:} A) increases as H+ concentration increases; B) decreases as blood pH decreases; C) decreases in resting muscle tissue; D) increases as blood pH decreases; E) increases as O$_{2}$ concentration decreases
\\ \textit{Note:} The correct answer is based on the Bohr effect, which states that an increase in carbon dioxide concentration and a decrease in blood pH lead to a reduction in hemoglobin's affinity for oxygen, facilitating oxygen release in actively respiring tissues.
\item \textbf{Answer:} A
\\ \textit{Options:} A) The temperature during embryonic development determines the sex.; B) The presence of an X or Y chromosome in the mother determines the sex of the offspring.; C) The presence of the X chromosome determines the sex.; D) The sex is determined by environmental factors after birth.; E) The presence of two Y chromosomes determines that an individual will be male.
\\ \textit{Note:} The correct answer highlights that in many species, including some reptiles, temperature can influence the expression of sex-determining genes during embryonic development, although in humans, sex is primarily determined genetically by the presence of the Y chromosome.
\item \textbf{Answer:} A
\\ \textit{Options:} A) They have flowers that bloom for several years before producing seeds; B) They have a unique leaf structure; C) They can survive in extremely cold climates; D) They produce large seeds; E) Their seeds can remain dormant for up to a century before germinating
\\ \textit{Note:} The correct answer highlights that cycads exhibit an extraordinary reproductive strategy where their flowers can persist for several years before seed production, distinguishing them from other seed plants that typically have shorter blooming periods.
\item \textbf{Answer:} A
\\ \textit{Options:} A) produce ATP; B) create carbon dioxide; C) make FADH 2; D) synthesize proteins; E) make glucose
\\ \textit{Note:} The energy released by electrons during their transfer through the electron transport chain is harnessed to pump protons across the mitochondrial membrane, creating a proton gradient that drives ATP synthase to produce ATP.
\end{enumerate}

\section*{True / False}
\begin{enumerate}[label=\arabic*.]
\setcounter{enumi}{5}
\item \textbf{Answer:} False
\\ \textit{Options:} A) True; B) False
\\ \textit{Note:} Laparoscopic antireflux surgery is a safe and effective treatment for GERD even in elderly patients, warranting low morbidity and mortality rates and a significant improvement of symptoms comparable to younger patients.
\item \textbf{Answer:} True
\\ \textit{Options:} A) True; B) False
\\ \textit{Note:} A greater association than hitherto acknowledged, between ascitis volume and anthropometric measurements from one side, and long-term rehospitalization and mortality from the other, was demonstrated in male stable alcoholic cirrhotics. Further studies with alcoholic and other modalities of cirrhosis including women are recommended.
\item \textbf{Answer:} False
\\ \textit{Options:} A) True; B) False
\\ \textit{Note:} The rabbit is a good model to be used in training of surgery, with a low morbi-mortality, able to be anesthetized intramuscularly, with no need of pre-operative fasting and does not present hypoglycemia even with the extended fasting period.
\item \textbf{Answer:} True
\\ \textit{Options:} A) True; B) False
\\ \textit{Note:} Postoperative numbness occurs in most patients receiving nasal microfat injections. Partial to complete recovery of nasal tip sensation can be expected to occur over a 3-month period.
\item \textbf{Answer:} True
\\ \textit{Options:} A) True; B) False
\\ \textit{Note:} A fetal anatomic survey on follow-up sonograms may identify unanticipated fetal anomalies, especially when the indication is for fetal growth.
\end{enumerate}

\section*{Long Form Response}
\begin{enumerate}[label=\arabic*.]
\setcounter{enumi}{10}
\item \textbf{Answer:} Prokaryotic and eukaryotic cells exhibit fundamental differences in their structures and functions that have significant implications for biological processes. Prokaryotic cells, which include bacteria and archaea, are characterized by their lack of a nucleus and membrane-bound organelles, with their genetic material organized in a single circular chromosome located in the nucleoid region. In contrast, eukaryotic cells, found in organisms such as plants, animals, and fungi, possess a defined nucleus containing linear chromosomes, as well as various organelles that compartmentalize cellular functions, enabling more complex metabolic processes. These structural differences influence cellular functions, including replication, gene expression, and energy production, highlighting the efficiency and adaptability of eukaryotic cells for multicellular organization and specialized functions. Additionally, the simplicity of prokaryotic cells allows for rapid reproduction and evolutionary adaptability, which plays a crucial role in ecological interactions and human health.
\\ \textit{Note:} The answer accurately delineates the structural differences between prokaryotic and eukaryotic cells, such as the presence of a nucleus and membrane-bound organelles in eukaryotes, and emphasizes how these differences affect their functions and broader biological processes.
\item \textbf{Answer:} Living organisms are defined by essential characteristics such as complex cellular organization, metabolic processes that enable energy consumption, movement, responsiveness to environmental stimuli, growth, reproduction, and the ability to adapt to changing conditions. These traits highlight the dynamic nature of life, which transcends mere chemical components. While the presence of the right elements and compounds is necessary for life, it is the intricate interactions and organizational structures among these components that give rise to the properties of living systems. Thus, life is not simply a collection of molecules but rather a sophisticated and integrated network of biochemical processes that result in the emergent properties we associate with living organisms.
\\ \textit{Note:} correct because it identifies that living organisms possess specific characteristics, such as cellular organization and metabolic processes, which demonstrate that life is a complex, dynamic system that cannot be reduced to merely having the right chemical components, as those components must interact in organized ways to support life functions.
\end{enumerate}

\vfill
\noindent\textit{End of Answer Key}
\end{document}